\hypertarget{native-533-lower}{}\label{native-533-lower}
\fancyhead[L]{\small KNITLMAYER ÜBER PLANETENABSTÄNDE.}
\begin{center}\small Göttingische gelehrte Anzeigen. Stück 179. S. 1777..1782. 10. November 1810.\\[3pt]\rule{.72\linewidth}{.3pt}\end{center}

\emph{Versuch einer genauen Darstellung des Progressionsverhältnisses der Planeten- und Trabantenabstände von ihren Zentralkörpern.} Von \textsc{Ferdinand Knitlmayer}, Hauptmann in der kaiserlich-königlichen österreichischen Armee. 1808. 48 Seiten im Oktavformat. Brünn.

Verschiedene Wahrheiten und Verhältnisse in der Astronomie sind auf dem Weg der Erfahrung oder des Probierens gleichsam erraten worden. Sie haben sich später vollständig bestätigt und zu wichtigen und weitreichenden Folgerungen geführt. Dieser glückliche Erfolg hat viele dazu ermuntert, die verschiedensten Zahlenverhältnisse in unserem Sonnensystem auf gut Glück miteinander zu vergleichen und einen Zusammenhang zwischen ihnen zu suchen. Hierher gehören vor allem die mittleren Abstände der Planeten von der Sonne. Soweit man die Planeten als einzelne, bereits gebildete Körper betrachtet, sind sowohl ihre Abstände vom Zentralkörper als auch alle übrigen Bestimmungsgrößen ihrer Bahnen tatsächlich als ganz unabhängig voneinander und willkürlich anzusehen. Es gibt überhaupt keinen Grund, Beziehungen zwischen ihnen vorauszusetzen. Nimmt man jedoch an, dass unser Sonnensystem seine ursprüngliche Bildung mechanischen Ursachen verdankt, so wäre durchaus denkbar, dass es zwischen den Elementen der einzelnen Planeten und etwa ihren Massen und Dichten Verhältnisse gibt. Könnte man solche Verhältnisse wirklich eindeutig erkennen, so ließe sich daraus mit etwas Glück vielleicht zurückschließen, wie unser Planetensystem ursprünglich entstanden ist. Darüber haben wir eigentlich nur vage und unzulässige Hypothesen. Bislang haben alle Bemühungen dieser Art, mit denen sich manche Liebhaber der Astronomie abgequält haben, zu gar nichts geführt. Kaum ein mathematischer Astronom wird darin mehr sehen als eine verzeihliche, aber unfruchtbare Spielerei.

In Deutschland hat man eine Zeit lang viel Gewicht auf ein vermeintliches Gesetz gelegt, in das man die Planetenabstände zwingen wollte. Man betrachtete sie als Summen aus einer Konstanten und den Gliedern einer geometrischen Progression, die sich von Planet zu Planet verdoppelt. So erhielt man keine wirkliche Übereinstimmung, sondern nur eine grobe Näherung. Merkur musste man ganz aus der Reihe herausnehmen, da sein Abstand nicht jener Summe, sondern der Konstanten selbst nahekam. Und für den Platz zwischen Mars und Jupiter haben sich sogar vier Anwärter gefunden. Vesta, die zuletzt bekannt wurde, ist bei Weitem die hellste und wahrscheinlich die größte von allen. Doch gerade sie weicht am stärksten von der Progression ab.

Die vorliegende kleine Schrift will ein anderes Gesetz aufstellen, nach dem sich die Planetenabstände nicht nur näherungsweise, sondern genau richten sollen. Herr \textsc{Knitlmayer} geht von einer geometrischen Progression aus. Er nimmt an, dass die wirklichen mittleren Abstände der einzelnen Planeten umso stärker von den Gliedern dieser Progression abweichen, je größer ihre Massen und je größer der Unterschied zwischen ihrer Dichte und der Dichte der Sonne sind. Die mittleren Abstände der verschiedenen Planeten sollen sich also durch die Formel
\SourceLeaf{534}
\[ a \,.\, 2^n + b\,m + c\,(D - \mathrm{1}) \]
darstellen lassen. Dabei bedeutet $m$ die Masse jedes Planeten, $D$ seine Dichte, wobei die der Sonne als Einheit genommen wird. $a$, $b$ und $c$ sind konstante Koeffizienten. $n$ bedeutet für Merkur 1, für Venus 2, für die Erde 3, für Mars 4, für die neuen Planeten 5, für Jupiter 6, für Saturn 7 und für Uranus 8.

Diese Hypothese hat an sich keinerlei Grundlage a priori. Um sie beim richtigen Namen zu nennen, ist sie aus der Luft gegriffen. Um sie zu prüfen, müsste man, sollte man meinen, die Werte der Konstanten $a$, $b$, $c$ durch Elimination aus den mittleren Entfernungen der drei Planeten bestimmen, deren Masse und Dichte am besten bekannt sind, also der Erde, des Jupiter und des Saturn. Dann müsste man die Formel an Planeten prüfen, deren Masse und Dichte einigermaßen genau bekannt sind, also an Venus und Uranus. Da wir die Dichten von Venus und Uranus eigentlich nur ungefähr kennen, müsste man sich auch dann noch hüten, einer etwaigen Übereinstimmung zu viel Gewicht beizumessen. Der Verfasser geht aber ganz anders vor. Er setzt sofort $a = \mathrm{0{,}075}$ an, ohne dass dafür ein anderer Grund ersichtlich wäre, als dass man dieselbe geozentrische Progression der oben erwähnten Hypothese zugrunde gelegt hatte. Die Koeffizienten $b$ und $c$ bestimmt er aus den mittleren Entfernungen der Erde und des Jupiter. Übertragen wir seine größtenteils sehr verworrene Darstellung in unsere eigene, so ergibt sich $b = \mathrm{427{,}435}$, $c = \mathrm{0{,}135143}$ (die Masse der Sonne $= \mathrm{1}$ gesetzt). Für Venus, Mars, Saturn und Uranus nimmt er dann die Massen an, die \textsc{Wurm} und \textsc{Laplace} angegeben haben. Daraus berechnet er mit seiner hypothetischen Formel die Dichten und daraus in Verbindung mit den Massen das Volumen und den scheinbaren Halbmesser in der Entfernung 1. Dass seine berechneten Werte recht gut mit den beobachteten übereinstimmen, sieht er als Beweis für die Richtigkeit seiner Hypothese an. Dann fällt es ihm nicht schwer, bei Merkur und den vier neuen Planeten die an sich unbekannten Massen und Dichten aus den von \textsc{Schröter} bestimmten scheinbaren Durchmessern abzuleiten. Auf den ersten Blick könnte diese angebliche Übereinstimmung manchen vielleicht blenden. Wir wollen daher nur angeben, wie der Verfasser bei seinen Rechnungen vorgeht. Bei Vesta, Saturn und Uranus ergibt die obige Formel für $D - \mathrm{1}$ einen negativen Wert $= -p$. Daher erklärt der Verfasser diese Planeten für weniger dicht als die Sonne. Doch anstatt nun die Dichte $= \mathrm{1} - p$ zu setzen, rechnet er so, als wäre sie $= -\dfrac{1}{1+p}$. Nach dieser Probe der mathematischen Kenntnisse des Verfassers wird man uns wohl alle weiteren Bemerkungen erlassen.

Der Verfasser wendet seine Hypothese nun auch auf die Trabantensysteme der Hauptplaneten an. Wenn wir seiner eigenen, zum Teil unverständlichen und widersprüchlichen Darstellung nicht folgen, sondern seine eigentliche Idee aus den Anwendungen erschließen, die er tatsächlich gibt, geht er folgendermaßen vor. Er nimmt an, dass sich die mittleren Abstände der einzelnen Trabanten ein und desselben Hauptplaneten ebenfalls durch die Formel
\[ \mathrm{0{,}075} \,.\, 2^n + \mathrm{427{,}435}\, m + \mathrm{0{,}135143}\,(D - \mathrm{1}) \]
darstellen lassen (wobei $m$ die Masse, $D$ die Dichte und $n$ die Ordnungszahl jedes Trabanten bedeutet). Nur soll man dann als Längeneinheit nicht wie oben die mittlere Entfernung der Erde von der Sonne nehmen, sondern für jeden Planeten eine andere. Sie beträgt nach der Vorschrift des Verfassers $\dfrac{A\BareRoot[3]{M}}{B}$, wobei $A$ den mittleren Abstand des Merkur von der Sonne bezeichnet, $B$ den mittleren Abstand des vom Trabanten begleiteten Hauptplaneten von der Sonne und $M$ dessen Masse (die Masse der Sonne stets $= \mathrm{1}$ gesetzt). Der Verfasser scheint diesen Wert der Längeneinheit so gewählt zu haben, dass sich damit eine leidliche \SourceLeaf{535}Übereinstimmung für den Trabanten der Erde ergibt. Dieser Wert hat aber keine weitere theoretische Grundlage. Die Trabanten der anderen Planeten lassen sich auch nicht zur Prüfung der Hypothese heranziehen, da wir ihre Dichten überhaupt nicht kennen. In der ganzen Anwendung der Hypothese auf die Trabantensysteme können wir deshalb nichts weiter erkennen als ein leeres Spiel mit Zahlen.
